% Part: computability
% Chapter: recursive-functions
% Section: general-recursive-functions

\documentclass[../../../include/open-logic-section]{subfiles}

\begin{document}

\olfileid{cmp}{rec}{gen}
\olsection{الدوال العودية العامة}

ولنا سبيل آخر إلى تعيين مجموعة من الدوال الكلية. نقول إن الدالة الكلية
$f(x,\vec z)$ \emph{منتظمة} متى كان لكل متتالية أعداد طبيعية
$\vec z$ عدد $x$ يحقق $f(x,\vec z)=0$. فالدوال المنتظمة، لا غير،
هي التي يعطي البحث غير المحدود فيها دالة
كلية. فلنقصر هذا البحث، احترازًا من الجزئية، على الدوال المنتظمة:

\begin{defn}
\ollabel{defn:general-recursive}
\emph{الدوال العودية العامة} تؤلف أصغر مجموعة من الدوال، بمختلف الرتب، من
الأعداد الطبيعية إلى الأعداد الطبيعية، تضم الصفر
واللاحق والإسقاطات، وتنغلق تحت التركيب والعودية البدائية، وتحت البحث
غير المحدود إذا أجري على الدوال \emph{المنتظمة}.
\end{defn}

فكل دالة عودية عامة كلية، كما يظهر من التعريف. وإنما يفترق
\olref{defn:general-recursive} عن \olref[par]{defn:recursive-fn} بأن
الثاني يجيز الدوال العودية الجزئية في أثناء البناء، ولا يشترط
الكلية إلا في الدالة التي ينتهي إليها. فكلمة «عامة»،
وهي تسمية موروثة من تاريخ البحث، قد توهم خلاف المقصود؛ إذ
\olref{defn:general-recursive} يبدو \emph{أقل} عمومًا من
\olref[par]{defn:recursive-fn}. غير أن هذا الفرق في ظاهر التعريف وحده:
فالدوال العودية العامة والدوال العودية، وإن اختلف طريق تعريفهما،
مجموعة واحدة.

\end{document}
